\chapter{Inledning}
Traditionellt har cybersäkerhet främst fokuserat på tekniska skyddsmekanismer, som brandväggar, kryptering och intrångsdetektion, men med digitaliseringens raska framsteg och förekomsten av nya angreppsmetoder har skyddet mot dessa hot blivit allt mer komplicerat \cite{cybersakerhet_kommuner}. Enligt \cite{heartfield_loukas_se, algarni2019se_survey} har hotaktörer börjat i allt större grad utnyttja mänskliga faktorer snarare än enbart tekniska sårbarheter. Social engineering definieras därmed som en angreppsform där psykologisk manipulation används för att få individer att lämna ut känslig information eller utföra handlingar som äventyrar säkerheten \cite{psychology_social_engineering}. Därför bör social engineering förstås som ett socio-tekniskt fenomen där tekniska system, organisatoriska strukturer och mänskligt beslutsfattande används för att begå brott \cite{washo2021interdisciplinary}.

Eftersom människan spelar en avgörande roll i social engineering-relaterade incidenter har den Europeiska unionens cybersäkerhetsbyrå (ENISA) börjat betona vikten av hur individuella beteenden, organisatoriska kulturer och arbetsprocesser påverkar förebyggandet av dessa attacker \cite{enisa2018cybersecurity_culture}. Detta innebär att endast tekniska skyddsåtgärder inte är tillräckliga för att hantera cyberhotet. Därtill har utvecklingen av artificiell intelligens (AI) ytterligare förändrat social engineering som ett cyberhot. Generativ AI har börjat användas för att producera mer avancerade attacker, vilket i sin tur har försvårat identifiering och detektion \cite{zhang2024digital_deception}. Samtidigt har AI också blivit ett verktyg för att stärka försvarssystemen genom avancerad mönsterigenkänning och hotanalys. Detta har lett till en kontinuerlig kapprustning där både hotaktörer och organisationer strävar efter att ha försprång.

\section{Syfte}
Syftet med denna kandidatavhandling är att analysera social engineering som ett cybersäkerhetshot med ett särskilt fokus på samspelet mellan tekniska och psykologiska faktorer. Avhandlingen ämnar att föra fram hur social engineering-attacker fungerar, varför de är effektiva och hur de kan förstås inom ramen för cybersäkerhetens socio-tekniska perspektiv. Som grund för avhandlingen används litteraturstudier där vetenskapliga artiklar, översiktsstudier och relevanta rapporter analyseras och sammanställs. Målet med avhandlingen är att bidra till en fördjupad förståelse av social engineering genom att identifiera centrala utmaningar kopplade till detektion och förebyggande av denna typ av cyberhot. 

\section{Forskningsfrågor}
I avhandlingen kommer fyra forskningsfrågor att behandlas. Inledningsvis förklaras social engineering i kontexten av modern cybersäkerhet. Därefter analyseras de psykologiska och tekniska mekanismer som möjliggör social engineering som ett effektivt cyberhot. Slutligen undersöks metoder som finns för att motverka, identifiera och filtrera social engineering i praktiken. Forskningsfrågorna ser ut som följande:
\begin{enumerate}
    \item Vilka psykologiska mekanismer bidrar till att social engineering-attacker är effektiva?
    \item Vilka tekniska metoder och attacktyper används i dagens social engineering-attacker?
    \item Vilka identifierings- och detektionsmetoder kan användas för att upptäcka social engineering?
    \item I vilken grad kan artificiell intelligens och automatiserade system bidra till att förebygga och motverka social engineering?
\end{enumerate}


\section{Disposition}